(1) 定义\(\hat N=a^\dagger a+(\sigma_z+I)/2\)。利用\([a^\dagger a,a]=-a\)与\([\sigma_z,\sigma_+]=2\sigma_+\)可见\([\hat N,a\sigma_+]=0\)，其Hermitian共轭亦然，故\([\hat N,H]=0\)。在\(\hat N=n+1\)子空间基底\((|e,n\rangle,|g,n+1\rangle)\)中，哈密顿量为\(\hbar\omega_c(n+\tfrac12)I+\frac\hbar2\begin{pmatrix}\Delta&2g\sqrt{n+1}\\2g\sqrt{n+1}&-\Delta\end{pmatrix}\)，\(\Delta=\omega_q-\omega_c\)。
(2) 迹零的非平凡矩阵平方是\([\Delta^2+4g^2(n+1)]I\)，所以能级\(E_{n,\pm}=\hbar\omega_c(n+\tfrac12)\pm\frac\hbar2\Omega_n\)，\(\Omega_n=\sqrt{\Delta^2+4g^2(n+1)}\)。共振\(\Delta=0\)时成对dressed态能差为\(2\hbar g\sqrt{n+1}\)；真空\(n=0\)也有非零\(2\hbar g\)劈裂。
(3) 在每个二能级子空间，去掉公共能量相位后演化算符为\(U=\cos(\Omega_nt/2)I-i\sin(\Omega_nt/2)(\Delta\sigma_z+2g\sqrt{n+1}\sigma_x)/\Omega_n\)。初始\((1,0)^T\)时跃迁振幅为\(-i[2g\sqrt{n+1}/\Omega_n]\sin(\Omega_nt/2)\)，概率\(P_{e,n\to g,n+1}(t)=4g^2(n+1)\sin^2(\Omega_nt/2)/[\Delta^2+4g^2(n+1)]\)。失谐增加使最大转移率低于一。
(4) 若\(|\Delta|\gg g\sqrt{n+1}\)，展开\(\Omega_n=|\Delta|[1+2g^2(n+1)/\Delta^2+O(g^4n^2/\Delta^4)]\)，dressed态获得大约\(\pm\hbar g^2(n+1)/|\Delta|\)的二阶频移，按分支连接原子裸态时符号随\(\Delta\)变化。共振外的该色散近似要求耦合相对失谐小；若\(g\)接近\(\omega_c\)、\(\omega_q\)或存在多能级跃迁，旋转波近似和二能级模型失效。
